% GENERATED FILE. Edit canonical lessons, then run npm run generate:books.
% canonical-source-hash: fnv1a64:4ec8f9b6d70ba46b
% canonical-lessons: MR-C469-kathor, MR-C469-nirdayi, MR-C469-dusta, MR-C469-ugra, MR-C469-kapti

\chapter{Harsh Natures}
\label{ch:mr-harsh-natures}
\hlchaptermodality{469}

\begin{chapteropening}
\textbf{By the end of this chapter:} \emph{I can describe a harsh, merciless, wicked, fierce or deceitful person.}

Complete the last lesson of chapter 469: I can describe a harsh, merciless, wicked, fierce or deceitful person.
\end{chapteropening}

\section[\texorpdfstring{\mr{कठोर} (kaṭhor)}{कठोर (kaṭhor)}]{\mr{कठोर} (kaṭhor) — harsh; stern}

\label{lesson:MR-C469-kathor}



\begin{quote}
Before the new one: say the Marathi for sensitive, then the Marathi for helpful to others.
\end{quote}

\subsection*{You'll want to know: \mr{कठोर}}
\textbf{\mr{कठोर}} — \emph{kaṭhor} — \textquotedblleft{}harsh; stern\textquotedblright{}.

\begin{quote}
\mr{ते} \mr{शिक्षक} \mr{खूप} \mr{कठोर} \mr{आहेत}.
\end{quote}

\begin{quote}
\emph{te śikṣak khūp kaṭhor āhet.}
\end{quote}

\begin{quote}
That teacher is very stern.
\end{quote}

Unchanging: \mr{कठोर} \mr{शब्द}, harsh words; \mr{कठोर} \mr{नियम}, strict rules.

\begin{cousinweb}
\textbf{\mr{कठोर}} (\emph{kaṭhor}) is a tatsama, taken from Sanskrit \emph{kaṭhora}, \textquotedblleft{}hard, stiff; cruel\textquotedblright{} (CDIAL 2651).
\end{cousinweb}

\begin{grammarlens}[title={What you've built}]
One more word for describing what you lived and felt.
\end{grammarlens}

\subsection*{Your turn}
\begin{itemize}
\raggedright
  \item \emph{Say it:} \emph{kaṭhor}
  \item \emph{Say it:} \emph{kaṭhor}, once more
  \item \emph{Say it:} \emph{te śikṣak khūp kaṭhor āhet.} — That teacher is very stern.
  \item \emph{Recall:} say the Marathi for sensitive, then the Marathi for helpful to others, then say \emph{kaṭhor} again
\end{itemize}

\begin{tcolorbox}[breakable,colback=teal!4,colframe=teal!35!black,title={Before you move on}]
What does \mr{कठोर} mean? (\textquotedblleft{}Harsh; stern\textquotedblright{}.) Say it once more.
\end{tcolorbox}
\section[\texorpdfstring{\mr{निर्दयी} (nirdayī)}{निर्दयी (nirdayī)}]{\mr{निर्दयी} (nirdayī) — merciless}

\label{lesson:MR-C469-nirdayi}



\begin{quote}
Before the new one: say the Marathi for helpful to others, then the Marathi for harsh.
\end{quote}

\subsection*{You'll want to know: \mr{निर्दयी}}
\textbf{\mr{निर्दयी}} — \emph{nirdayī} — \textquotedblleft{}merciless\textquotedblright{}.

\begin{quote}
\mr{कुत्र्याला} \mr{मारणारा} \mr{माणूस} \mr{निर्दयी} \mr{असतो}.
\end{quote}

\begin{quote}
\emph{kutryālā mārṇārā māṇūs nirdayī asto.}
\end{quote}

\begin{quote}
A man who beats a dog is merciless.
\end{quote}

\mr{निर्}, without, + \mr{दया}, pity, + the ending -\mr{ई}.

\begin{cousinweb}
\textbf{\mr{निर्दयी}} (\emph{nirdayī}) is built on the tatsama \emph{nirdaya}, Sanskrit for \textquotedblleft{}pitiless, cruel\textquotedblright{} (Monier-Williams), with the Marathi adjective ending \emph{-ī}.
\end{cousinweb}

\begin{grammarlens}[title={What you've built}]
One more word for describing what you lived and felt.
\end{grammarlens}

\subsection*{Your turn}
\begin{itemize}
\raggedright
  \item \emph{Say it:} \emph{nirdayī}
  \item \emph{Say it:} \emph{nirdayī}, once more
  \item \emph{Say it:} \emph{kutryālā mārṇārā māṇūs nirdayī asto.} — A man who beats a dog is merciless.
  \item \emph{Recall:} say the Marathi for helpful to others, then the Marathi for harsh, then say \emph{nirdayī} again
\end{itemize}

\begin{tcolorbox}[breakable,colback=teal!4,colframe=teal!35!black,title={Before you move on}]
What does \mr{निर्दयी} mean? (\textquotedblleft{}Merciless\textquotedblright{}.) Say it once more.
\end{tcolorbox}
\section[\texorpdfstring{\mr{दुष्ट} (duṣṭa)}{दुष्ट (duṣṭa)}]{\mr{दुष्ट} (duṣṭa) — wicked}

\label{lesson:MR-C469-dusta}



\begin{quote}
Before the new one: say the Marathi for harsh, then the Marathi for merciless.
\end{quote}

\subsection*{You'll want to know: \mr{दुष्ट}}
\textbf{\mr{दुष्ट}} — \emph{duṣṭa} — \textquotedblleft{}wicked\textquotedblright{}.

\begin{quote}
\mr{गोष्टीतला} \mr{दुष्ट} \mr{राजा} \mr{शेवटी} \mr{हरला}.
\end{quote}

\begin{quote}
\emph{goṣṭītlā duṣṭa rājā śevṭī harlā.}
\end{quote}

\begin{quote}
The wicked king in the story was defeated in the end.
\end{quote}

Unchanging; also a noun, a villain.

\begin{cousinweb}
\textbf{\mr{दुष्ट}} (\emph{duṣṭa}) is a tatsama, taken from Sanskrit \emph{duṣṭa}, \textquotedblleft{}defiled, bad; a villain\textquotedblright{} (CDIAL 6472).
\end{cousinweb}

\begin{grammarlens}[title={What you've built}]
One more word for describing what you lived and felt.
\end{grammarlens}

\subsection*{Your turn}
\begin{itemize}
\raggedright
  \item \emph{Say it:} \emph{duṣṭa}
  \item \emph{Say it:} \emph{duṣṭa}, once more
  \item \emph{Say it:} \emph{goṣṭītlā duṣṭa rājā śevṭī harlā.} — The wicked king in the story was defeated in the end.
  \item \emph{Recall:} say the Marathi for harsh, then the Marathi for merciless, then say \emph{duṣṭa} again
\end{itemize}

\begin{tcolorbox}[breakable,colback=teal!4,colframe=teal!35!black,title={Before you move on}]
What does \mr{दुष्ट} mean? (\textquotedblleft{}Wicked\textquotedblright{}.) Say it once more.
\end{tcolorbox}
\section[\texorpdfstring{\mr{उग्र} (ugra)}{उग्र (ugra)}]{\mr{उग्र} (ugra) — fierce; sharp}

\label{lesson:MR-C469-ugra}



\begin{quote}
Before the new one: say the Marathi for merciless, then the Marathi for wicked.
\end{quote}

\subsection*{You'll want to know: \mr{उग्र}}
\textbf{\mr{उग्र}} — \emph{ugra} — \textquotedblleft{}fierce; sharp\textquotedblright{}.

\begin{quote}
\mr{रागावल्यावर} \mr{त्याचा} \mr{चेहरा} \mr{उग्र} \mr{दिसतो}.
\end{quote}

\begin{quote}
\emph{rāgāvlyāvar tyācā cehrā ugra disto.}
\end{quote}

\begin{quote}
When he gets angry his face looks fierce.
\end{quote}

Unchanging; of a face, a smell or a medicine: \mr{उग्र} \mr{वास}, a sharp smell.

\begin{cousinweb}
\textbf{\mr{उग्र}} (\emph{ugra}) is a tatsama, taken from Sanskrit \emph{ugra}, \textquotedblleft{}strong, terrible, fierce\textquotedblright{} (CDIAL 1631).
\end{cousinweb}

\begin{grammarlens}[title={What you've built}]
One more word for describing what you lived and felt.
\end{grammarlens}

\subsection*{Your turn}
\begin{itemize}
\raggedright
  \item \emph{Say it:} \emph{ugra}
  \item \emph{Say it:} \emph{ugra}, once more
  \item \emph{Say it:} \emph{rāgāvlyāvar tyācā cehrā ugra disto.} — When he gets angry his face looks fierce.
  \item \emph{Recall:} say the Marathi for merciless, then the Marathi for wicked, then say \emph{ugra} again
\end{itemize}

\begin{tcolorbox}[breakable,colback=teal!4,colframe=teal!35!black,title={Before you move on}]
What does \mr{उग्र} mean? (\textquotedblleft{}Fierce; sharp\textquotedblright{}.) Say it once more.
\end{tcolorbox}
\section[\texorpdfstring{\mr{कपटी} (kapṭī)}{कपटी (kapṭī)}]{\mr{कपटी} (kapṭī) — deceitful}

\label{lesson:MR-C469-kapti}



\begin{quote}
Before the new one: say the Marathi for wicked, then the Marathi for fierce.
\end{quote}

\subsection*{You'll want to know: \mr{कपटी}}
\textbf{\mr{कपटी}} — \emph{kapṭī} — \textquotedblleft{}deceitful\textquotedblright{}.

\begin{quote}
\mr{त्या} \mr{कपटी} \mr{माणसावर} \mr{विश्वास} \mr{ठेवू} \mr{नये}.
\end{quote}

\begin{quote}
\emph{tyā kapṭī māṇsāvar viśvās ṭhevū naye.}
\end{quote}

\begin{quote}
One should not trust that deceitful man.
\end{quote}

From \mr{कपट}, deceit. Unchanging.

\begin{cousinweb}
\textbf{\mr{कपटी}} (\emph{kapṭī}) is a tatsama, taken from Sanskrit \emph{kapaṭin}, \textquotedblleft{}deceitful\textquotedblright{} (Monier-Williams), built on \emph{kapaṭa}, \textquotedblleft{}fraud\textquotedblright{}.
\end{cousinweb}

\begin{grammarlens}[title={What you've built}]
One more word for describing what you lived and felt.
\end{grammarlens}

\subsection*{Your turn}
\begin{itemize}
\raggedright
  \item \emph{Say it:} \emph{kapṭī}
  \item \emph{Say it:} \emph{kapṭī}, once more
  \item \emph{Say it:} \emph{tyā kapṭī māṇsāvar viśvās ṭhevū naye.} — One should not trust that deceitful man.
  \item \emph{Recall:} say the Marathi for wicked, then the Marathi for fierce, then say \emph{kapṭī} again
\end{itemize}

\begin{tcolorbox}[breakable,colback=teal!4,colframe=teal!35!black,title={Before you move on}]
What does \mr{कपटी} mean? (\textquotedblleft{}Deceitful\textquotedblright{}.) Say it once more.
\end{tcolorbox}
